
\[
  Q_{2,2}(X_1,X_2,X_3,X_4)=X_1X_3-X_4.
\]

Indiquons pour finir l'origine des formules (2.3.1) et (2.3.3). On considère l'exemple 2.2.b), où les opérations \(\lambda^i\) sont définies par les puissances extérieures des \(A\)-Modules. On souhaite que \(\lambda_t:K\longrightarrow\Ghat(K)\) soit un \(\lambda\)-homomorphisme~; on doit donc définir sur \(\Ghat\) des lois telles que pour tous \(A\)-Modules \(\underline{E},\underline{F}\), on ait~:
\[
  \lambda_t(EF)=\lambda_t(E)\circ\lambda_t(F)~;\qquad
  \lambda_u(\lambda_t(E))=\Ghat^u(\lambda_t)(\lambda_u(E))~;
\]
or les relations (2.3.1) et (2.3.3) sont précisément les relations précédentes appliquées au cas où \(\underline{E}\) et \(\underline{F}\) sont des modules inversibles.

Nous verrons dans Exp. VI dans quelle mesure la propriété souhaitée est vérifiée.

\textbf{Définition 2.4.} \uline{Un pré}-\(\lambda\)-\uline{anneau} \(K\) \uline{sera appelé un} \(\lambda\)-\uline{anneau si l'homomorphisme de groupes abéliens} \(\lambda_t:K\longrightarrow1+K[[t]]^+\) \uline{est un} \(\lambda\)-\uline{homomorphisme} \uline{pour la structure de pré}-\(\lambda\)-\uline{anneau sur le deuxième membre définie en} 2.3.

Cela signifie donc qu'on a les relations~: \(\forall x,y\in K\),
\begin{equation}\tag{2.4.1}
  \lambda_t(1)=1+t~;\quad
  \lambda_t(xy)=\lambda_t(x)\circ\lambda_t(y)~;\quad
  \lambda_u(\lambda_t(x))=\Ghat(\lambda_t)\bigl(\lambda_u(x)\bigr).
\end{equation}

En termes des applications \(\lambda^i\), les deux dernières relations équivalent à : \(x,y\in K\), \(i.j\in\mathbb N\),
\begin{align}
  \lambda^i(xy)
  &=P_i(\lambda^1(x),\ldots,\lambda^i(x);
       \lambda^1(y),\ldots,\lambda^i(y)) \tag{2.4.2}\\
  \lambda^j(\lambda^i(x))
  &=Q_{i,j}(\lambda^1(x),\ldots,\lambda^{ij}(x)) \notag
\end{align}

